% Part: first-order-logic
% Chapter: completeness
% Section: outline

\documentclass[../../../include/open-logic-section]{subfiles}

\begin{document}

\iftag{FOL}
      {\olfileid{fol}{com}{out}}
      {\olfileid{pl}{com}{out}}

\olsection{طرح کلی برهان}

برهان قضیهٔ تمامیت قدری پیچیده است و در نخستین مطالعه به‌آسانی می‌توان
در آن راه را گم کرد. پس نخست طرحی کلی از برهان به دست می‌دهیم. گام نخست
تغییر دیدگاهی است که راهی به سوی برهان پیش چشممان می‌گذارد. اگر تمامیت
را به صورت «هرگاه $\Gamma \Entails !A$، آنگاه $\Gamma \Proves !A$» در
نظر بگیریم، حتی یافتن ایده‌ای برای برهان ممکن است دشوار باشد: برای نشان
دادن $\Gamma \Proves !A$ باید !!a{derivation} بیابیم، و به نظر نمی‌رسد
فرض $\Gamma \Entails !A$ در این کار هیچ کمکی بکند. در برخی دستگاه‌های
برهان می‌توان !!a{derivation} را مستقیماً ساخت، اما ما رویکردی اندکی
متفاوت در پیش می‌گیریم. تغییر دیدگاه لازم چنین است: تمامیت را می‌توان
این‌گونه نیز صورت‌بندی کرد: «اگر $\Gamma$ سازگار باشد، ارضاپذیر است.»
شاید بتوانیم اطلاعات موجود در~$\Gamma$ را همراه با فرض سازگاری آن به کار
بگیریم تا \iftag{FOL}{!!a{structure}}{!!a{valuation}} بسازیم که هر
\iftag{FOL}{!!{sentence}}{!!{formula}} موجود در~$\Gamma$ را ارضا کند.
سرانجام می‌دانیم در پی چه نوع
\iftag{FOL}{!!{structure}}{!!{valuation}}ای هستیم: چیزی که درست همان‌گونه
باشد که $\Gamma$ توصیفش می‌کند!{}

اگر $\Gamma$ فقط شامل \iftag{FOL}{!!{sentence}s
  اتمی}{!!{propositional variable}s} باشد، ساختن مدلی برای آن آسان
است.\iftag{FOL}{ فرض کنید همهٔ !!{sentence}s اتمی به صورت
  $\Atom{P}{a_1,\dots,a_n}$ باشند، که در آن $a_i$ها !!{constant}
  هستند.}{} کافی است
\iftag{FOL}{%
  !!a{domain}~$\Domain{M}$ و تعبیری برای~$P$ بیابیم چنان‌که
  $\Sat{M}{\Atom{P}{a_1,\ldots,a_n}}$. این کار چندان دشوار نیست:
  $\Domain{M} = \Nat$، $\Assign{\Obj c_i}{M} = i$، و برای هر
  $\Atom{P}{a_1,\ldots,a_n} \in \Gamma$، چندتاییِ $\tuple{k_1,
    \dots, k_n}$ را در $\Assign{P}{M}$ قرار می‌دهیم، که در آن $k_i$
  اندیس نماد ثابت~$a_i$ است (یعنی $a_i \ident \Obj c_{k_i}$).}{%
  !!a{valuation}~$\pAssign{v}$ چنان‌که $\pSat{v}{p}$ برای همهٔ $p \in
  \Gamma$. کافی است $\pAssign{v}(p) = \True$ هرگاه و تنها هرگاه
  $p \in \Gamma$.}

اکنون فرض کنید $\Gamma$ شامل !!a{formula}~$\lnot !B$ باشد که $!B$ اتمی
است. ممکن است نگران باشیم که ساخت
$\iftag{FOL}{\Struct{M}}{\pAssign{v}}$ با امکان صادق‌ساختن $\lnot !B$
تداخل کند. اما اینجا سازگاری~$\Gamma$ به کار می‌آید: اگر
$\lnot !B \in \Gamma$، آنگاه $!B \notin \Gamma$، وگرنه $\Gamma$
ناسازگار می‌بود. و اگر $!B \notin \Gamma$، آنگاه بنا بر ساخت ما از
$\iftag{FOL}{\Struct{M}}{\pAssign{v}}$،
$\iftag{FOL}{\Sat/{M}}{\pSat/{v}}{!B}$؛ بنابراین
$\iftag{FOL}{\Sat{M}}{\pSat{v}}{\lnot !B}$. تا اینجا همه‌چیز درست پیش
می‌رود.

اگر $\Gamma$ شامل فرمول‌های مرکب و نااتمی باشد چه؟ فرض کنید شامل
$!A \land !B$ باشد. برای صادق‌ساختن آن باید چنان عمل کنیم که گویی هر دو
$!A$ و~$!B$ در~$\Gamma$ هستند. و اگر $!A \lor !B \in \Gamma$، باید
دست‌کم یکی از آن‌ها را صادق سازیم؛ یعنی چنان عمل کنیم که گویی یکی از آن
دو در~$\Gamma$ است.

این ملاحظات ایدهٔ زیر را پیش می‌نهند: !!{formula}s دیگری به~$\Gamma$
می‌افزاییم تا (الف)~مجموعهٔ حاصل سازگار بماند؛ (ب)~برای هر
\iftag{FOL}{!!{sentence}}{!!{formula}} ممکن~$!A$، یا $!A$ در مجموعهٔ
حاصل باشد یا $\lnot !A$؛ و (ج)~هرگاه $!A \land !B$ در مجموعه باشد، هر
دو $!A$ و~$!B$ نیز در آن باشند، و اگر $!A \lor !B$ در مجموعه باشد،
دست‌کم یکی از $!A$ و~$!B$ نیز در آن باشد، و جز این‌ها. این کار را
(احتمالاً تا بی‌نهایت) ادامه می‌دهیم. مجموعهٔ همهٔ !!{formula}s افزوده
را~$\Gamma^*$ می‌نامیم. آنگاه ساخت بالا
\iftag{FOL}{!!a{structure}~$\Struct{M}$}{!!a{valuation}~$\pAssign{v}$}
به ما می‌دهد که می‌توانیم با استقرا اثبات کنیم همهٔ جمله‌های
در~$\Gamma^*$، و در نتیجه همهٔ جمله‌های در~$\Gamma$ را نیز ارضا می‌کند،
زیرا $\Gamma \subseteq \Gamma^*$. معلوم می‌شود تضمین (الف) و~(ب) کافی
است. مجموعه‌ای از جمله‌ها که (ب) برای آن برقرار باشد \emph{کامل} نامیده
می‌شود. پس وظیفهٔ ما گسترش مجموعهٔ سازگار~$\Gamma$ به مجموعهٔ سازگار و
کامل~$\Gamma^*$ خواهد بود.

\iftag{FOL}{%
این طرح یک دشواری دارد: اگر $\lexists[x][!A(x)] \in \Gamma$، امیدواریم
بتوانیم !!a{constant}~$c$ برگزینیم و در این فرایند $!A(c)$ را بیفزاییم.
اما از کجا بدانیم همیشه می‌توان چنین کرد؟ شاید زبان ما فقط چند
!!{constant} داشته باشد و برای هر یک از آن‌ها
$\lnot !A(c) \in \Gamma$. نمی‌توانیم $!A(c)$ را نیز بیفزاییم، زیرا
مجموعه ناسازگار می‌شود و دیگر نمی‌دانیم $\Struct{M}$ باید $!A(c)$ را
صادق سازد یا $\lnot !A(c)$ را. حتی ممکن است $\Gamma$ فقط شامل جمله‌های
زبانی باشد که هیچ نماد ثابتی ندارد (برای نمونه، زبان نظریهٔ مجموعه‌ها).

راه‌حل این مشکل آن است که در آغاز بی‌نهایت ثابت به زبان بیفزاییم و
جمله‌هایی نیز اضافه کنیم که آن‌ها را به‌شیوهٔ مناسب با سورها پیوند
می‌دهند. (البته باید بررسی کنیم که این کار نمی‌تواند ناسازگاری پدید
آورد.)

ساخت آغازین ما هنگامی خوب عمل می‌کند که در جمله‌های اتمی فقط
!!{constant}s داشته باشیم. اما زبان ممکن است شامل !!{function}s نیز
باشد. در این صورت یافتن توابع مناسب روی~$\Nat$ برای نسبت‌دادن به این
!!{function}s و درست‌کردن همه‌چیز ممکن است دشوار باشد. پس ترفند دیگری
به کار می‌بریم: به‌جای استفاده از $i$ برای تعبیر $\Obj c_i$، خود مجموعهٔ
!!{constant}s را دامنه می‌گیریم. آنگاه $\Struct M$ می‌تواند هر
!!{constant} را به خودش نسبت دهد: $\Assign{\Obj c_i}{M} = \Obj c_i$.
اما چرا تا پایان این راه نرویم؟ بگذارید $\Domain{M}$ مجموعهٔ همهٔ
\emph{ترم‌های} زبان باشد!{} در این صورت، نسبت‌دادن توابع (که ترم‌ها را
به‌عنوان ورودی می‌گیرند و ترم‌ها را به‌عنوان مقدار می‌دهند) به
!!{function}s روشن است: به !!{function}~$\Obj f^n_i$ تابعی را نسبت
می‌دهیم که با گرفتن $n$ ترم $t_1$، \dots،~$t_n$ به‌عنوان ورودی، ترم
$\Obj f^n_i(t_1, \dots, t_n)$ را به‌عنوان مقدار پدید می‌آورد.

آخرین قطعهٔ معما این است که با~$\eq$ چه کنیم. !!{predicate}~$\eq$ تعبیر
ثابتی دارد: $\Sat{M}{\eq[t][t']}$ هرگاه و تنها هرگاه
$\Value{t}{M} = \Value{t'}{M}$. حال اگر امور را چنان ترتیب دهیم که
!!{value} ترم~$t$ خود~$t$ باشد، این !!{structure} هیچ جمله‌ای به صورت
$\eq[t][t']$ را صادق نخواهد ساخت، مگر آنکه $t$ و~$t'$ عیناً یک ترم
باشند. البته این مشکل‌ساز است، زیرا تقریباً هر نظریهٔ جالب در زبانی با
!!{function}s، جمله‌هایی به صورت $\eq[t][t']$ را قضیه دارد که در آن‌ها
$t$ و~$t'$ یک ترم نیستند (برای نمونه، در نظریه‌های حساب:
$\eq[(\Obj 0+ \Obj 0)][\Obj 0]$). برای حل این مشکل دامنهٔ~$\Struct M$
را تغییر می‌دهیم: به‌جای استفاده از ترم‌ها به‌عنوان اشیای~$\Domain{M}$،
مجموعه‌هایی از ترم‌ها را به کار می‌بریم؛ هر مجموعه همهٔ ترم‌هایی را دربر
می‌گیرد که جمله‌های در~$\Gamma$ برابر‌بودنشان را ایجاب می‌کنند. برای
نمونه، اگر $\Gamma$ نظریه‌ای دربارهٔ حساب باشد، یکی از این مجموعه‌ها
شامل $\Obj 0$، $(\Obj 0 + \Obj 0)$، $(\Obj 0 \times \Obj 0)$ و جز
این‌ها خواهد بود. این همان مجموعه‌ای است که به $\Obj 0$ نسبت می‌دهیم،
و معلوم خواهد شد همین مجموعه مقدار همهٔ ترم‌های درون آن، از جمله
$(\Obj 0 + \Obj 0)$، نیز هست. بنابراین جملهٔ
$\eq[(\Obj 0+ \Obj 0)][\Obj 0]$ در این !!{structure} بازنگری‌شده صادق
خواهد بود.}{}

پس کارمان چنین خواهد بود. نخست خواص مجموعه‌های سازگار !!{complete} را
بررسی می‌کنیم؛ به‌ویژه اثبات می‌کنیم که !!a{complete} مجموعهٔ سازگار
شامل $!A \land !B$ است هرگاه و تنها هرگاه شامل هر دو $!A$ و~$!B$ باشد،
شامل $!A \lor !B$ است هرگاه و تنها هرگاه شامل دست‌کم یکی از آن‌ها باشد،
و جز این‌ها (\olref[ccs]{prop:ccs}).\iftag{FOL}{ سپس مجموعه‌های جملهٔ
  «اشباع‌شده» را تعریف و بررسی می‌کنیم. مجموعهٔ اشباع‌شده مجموعه‌ای است
  که شرطی‌هایی را دربر دارد که هر !!{sentence} مسور را به نمونه‌هایی از
  آن پیوند می‌دهند (\olref[hen]{defn:henkin-exp}). نشان می‌دهیم که هر
  مجموعهٔ سازگار~$\Gamma$ را همواره می‌توان به مجموعهٔ
  اشباع‌شدهٔ~$\Gamma'$ گسترش داد (\olref[hen]{lem:henkin}). اگر مجموعه‌ای
  سازگار، اشباع‌شده و !!{complete} باشد، این خاصیت را نیز دارد که شامل
  \iftag{prvEx}{$\lexists[x][!A(x)]$ است هرگاه و تنها هرگاه برای یک ترم
    بستهٔ~$t$ شامل $!A(t)$ باشد}{}\iftag{defEx,defAll}{}{ و
  }\iftag{prvAll}{$\lforall[x][!A(x)]$ است هرگاه و تنها هرگاه برای همهٔ
    ترم‌های بستهٔ~$t$ شامل $!A(t)$ باشد}{}
  (\olref[hen]{prop:saturated-instances}).}{}
سپس مجموعهٔ سازگار~$\iftag{FOL}{\Gamma'}{\Gamma}$ را که
\iftag{FOL}{اشباع‌شده نیز هست}{} می‌گیریم و نشان می‌دهیم می‌توان آن را
به مجموعهٔ~$\Gamma^*$ گسترش داد که
\iftag{FOL}{اشباع‌شده، سازگار و}{سازگار و} !!{complete} است
(\olref[lin]{lem:lindenbaum}). از همین مجموعهٔ $\Gamma^*$ برای تعریف
\iftag{FOL}{مدل ترمی~$\Struct M(\Gamma^*)$}{!!{valuation}~$\pAssign
  v(\Gamma^*)$} استفاده خواهیم کرد. \iftag{FOL}{مدل ترمی مجموعهٔ ترم‌های
  بسته را به‌عنوان دامنه دارد و تعبیر !!{predicate}s آن به‌وسیلهٔ
  !!{sentence}s اتمی}{ارزش‌گذاری به‌وسیلهٔ !!{propositional variable}s}
در~$\Gamma^*$ تعیین می‌شود (\olref[mod]{defn:termmodel}). با بهره‌گیری از
خواص مجموعه‌های سازگارِ کامل\iftag{FOL}{ و اشباع‌شده}{} نشان خواهیم داد
که واقعاً
$\iftag{FOL}{\Sat{M(\Gamma^*)}}{\pSat{v(\Gamma^*)}}{!A}$ هرگاه و تنها
هرگاه $!A \in \Gamma^*$ (\olref[mod]{lem:truth})، و بنابراین به‌ویژه
$\iftag{FOL}{\Sat{M(\Gamma^*)}}{\pSat{v(\Gamma^*)}}{\Gamma}$.
\iftag{FOL}{در پایان، بررسی می‌کنیم که اگر $\Gamma$ شامل~$\eq$ نیز باشد
  مدل ترمی را چگونه تعریف کنیم (\olref[ide]{defn:term-model-factor}) و
  نشان می‌دهیم که این مدل~$\Gamma^*$ را ارضا می‌کند
  (\olref[ide]{lem:truth}).}{}

\end{document}
